Let $A$ be a Macaulay ring and $\mathfrak{a}$ an ideal in $A$ generated by a prime sequence $\{a_1, \cdots, a_j\}$. We have $\mathfrak{a}^n: Aa_j = \mathfrak{a}^{n-1}$ for every $n$.

Let $x$ be an element of $A$ such that $xa_j \in \mathfrak{a}^n$. We set $\mathfrak{b} = Aa_1 + \cdots + Aa_{j-1}$. Since $\mathfrak{a}^n = \mathfrak{b}^n + \mathfrak{a}^{n-1} a_j$, there exists an element $y$ of $\mathfrak{a}^{n-1}$ such that $(x-y)a_j \in \mathfrak{b}^n$. As $a_j$ is prime to $\mathfrak{b}$, it is prime to $\mathfrak{b}^n$ (Lemma 5), whence $x-y \in \mathfrak{b}^n$. Therefore $x \in \mathfrak{a}^{n-1}$, and we have proved the inclusion $\mathfrak{a}^n: Aa_j \subset \mathfrak{a}^{n-1}$. Since the opposite

ideal, and there exist elements $b_i$ of $\mathfrak{q}^{s-1}$ ($i \in I, i \neq k$) such that $\bar{z}_k = \sum_{i \in I, i \neq k} b_i \bar{a}_i$, i.e., such that $z_k - \sum_{i \in I} b_i a_i$ is an element $z'_k$ of $\mathfrak{q}^{s+1}$. Setting $z'_i = z_i + b_i a_k$ for $i \in I, i \neq k$, and $z'_i = z_i$ for $i \notin I$, we get a relation $\sum_{i=1}^{d} z'_i a_i = 0$ in which $v(z'_i) \geq v(z_i)$ for every $i$ and $v(z'_k) > v(z_k)$.

Now, among the relations $ya_j = \sum_{i=1}^{j-1} x_i a_i$ ($x_i \in A$) we choose one which has the following two properties: (a) $\min_i (v(x_i))$ has the greatest possible value, say $s$; (b) the number of indices $i$ such that $v(x_i) = s$ is the smallest possible. Then we have $s = v(y)$. In fact $s > v(y)$ is obviously impossible. On the other hand, if $s < v(y)$, we transform, as above, the relation $ya_j - \sum_{i=1}^{j-1} x_i a_i = 0$: the coefficient $y$ of $a_j$ is then unchanged, whereas, either $s$ is increased, or the number of indices $i$ such that $v(x_i) = s$ is decreased. This is impossible. Thus $v(y) = \min_i (v(x_i))$. Transforming, as above, the relation $ya_j - \sum_{i=1}^{j-1} x_i a_i = 0$, this time with $ya_j$ playing the part of $z_k a_k$, we get a relation $y_1 a_j - \sum_{i=1}^{j-1} x'_i